\documentclass[11pt]{amsart}

\usepackage[T1]{fontenc}
\usepackage{lmodern}
\usepackage[final]{microtype}
\usepackage{amsmath,amssymb,mathtools}
\usepackage[
  letterpaper,
  top=1.05in,
  bottom=1.08in,
  inner=1.12in,
  outer=1.12in,
  headheight=14pt,
  headsep=18pt
]{geometry}
\usepackage{xcolor}
\definecolor{AMSBlue}{HTML}{174A72}
\definecolor{AMSCite}{HTML}{7A3E20}
\usepackage[
  colorlinks=true,
  linkcolor=AMSBlue,
  citecolor=AMSCite,
  urlcolor=AMSBlue,
  pdfborder={0 0 0},
  breaklinks=true
]{hyperref}

\hypersetup{
  pdftitle={Positivity of the Quartic Solution of a Higher Discrete Painleve Equation},
  pdfsubject={Quartic discrete Painleve recurrence and quantum minimal surfaces},
  pdfkeywords={discrete Painleve equation, quartic potential, orthogonal polynomials, positivity}
}
\urlstyle{same}

\linespread{1.025}
\setlength{\emergencystretch}{2em}
\allowdisplaybreaks[2]
\raggedbottom

\numberwithin{equation}{section}

\newtheorem{theorem}{Theorem}[section]
\newtheorem{proposition}[theorem]{Proposition}
\newtheorem{lemma}[theorem]{Lemma}
\newtheorem{corollary}[theorem]{Corollary}
\theoremstyle{remark}
\newtheorem{remark}[theorem]{Remark}

\newcommand{\C}{\mathbb C}
\newcommand{\N}{\mathbb N}
\newcommand{\dd}{\,\mathrm d}
\newcommand{\ii}{\mathrm i}
\newcommand{\sgn}{\operatorname{sgn}}
\newcommand{\Span}{\operatorname{span}}
\newcommand{\e}{\mathrm e}
\newcommand{\fall}[1]{^{\underline{#1}}}

\title[Positivity of the quartic solution]
{Positivity of the Quartic Solution of a Higher\\
Discrete Painlev\'e Equation}

% Add author, affiliation, and contact information before submission.
\author{}
\date{}

\subjclass[2020]{Primary 34M55; Secondary 33C45, 15B52, 53A10}
\keywords{discrete Painlev\'e equation, indefinite orthogonal polynomials,
quartic potential, total positivity, quantum minimal surface}

\begin{document}

\begin{abstract}
We prove positivity of the solution of the quartic discrete
Painlev\'e-type recurrence arising from the quantum minimal surface
\(w_2=w_1^3\).  The initial data are given by a product of Bessel
functions of orders \(\pm\frac14\).  The proof converts the Bessel
amplitude into a positive Laplace integral and introduces two auxiliary
quartic source variables.  The four residue classes modulo \(4\) then
reduce to a single family of bi-moment matrices.  A triangular
factorization, followed by Andr\'eief's identity, fixes the strict sign
of every relevant principal minor.  Consequently all signed
orthogonal-polynomial norms are positive, and their consecutive ratios
give \(v_n(a)>0\) for every \(a>0\) and \(n\geq0\).
\end{abstract}

\maketitle

\section{Introduction}

Discrete Painlev\'e-type recurrences occur in matrix regularizations of
minimal surfaces and in the recurrence coefficients of orthogonal
polynomials.  For the quantum minimal surface associated with the
holomorphic curve \(w_2=w_1^{d-1}\), Felder and Hoppe
\cite{FelderHoppe2025} constructed an indefinite Hermitian form whose
orthogonal-polynomial norm ratios obey
\[
V_n+\frac1{a^2}\sum_{j=0}^{d-2}
\prod_{k=0}^{d-2}V_{n+j-k}=\frac{n+1}{a}.
\]
This connects the discrete equation with the operator formulation of
quantum minimal surfaces originating in \cite{CornalbaTaylor1999}; see
also \cite{ClarksonEtAl2025} for special solutions in the cubic case.

For \(d=4\), after setting \(v_n=V_n/a\) and
\(\epsilon=a^{-2}\), the recurrence becomes
\begin{equation}
v_n\left(
1+v_{n-1}v_{n-2}
+v_{n+1}v_{n-1}
+v_{n+2}v_{n+1}
\right)
=(n+1)\epsilon ,
\qquad n\geq0,
\label{eq:quartic-recurrence}
\end{equation}
where \(v_n=0\) for \(n<0\).  The two initial values are determined by
\begin{equation}
v_0=-\frac1a\frac{h'(a)}{h(a)},
\qquad
v_1=-\frac1a\frac{h''(a)}{h'(a)},
\label{eq:initial-values}
\end{equation}
with
\begin{equation}
\begin{split}
h(a)=\frac{\pi^2a}{4}\bigg[
&J_{-1/4}\left(\frac{a^2}{4}\right)^2
-\sqrt2\,J_{-1/4}\left(\frac{a^2}{4}\right)
J_{1/4}\left(\frac{a^2}{4}\right)\\
&+J_{1/4}\left(\frac{a^2}{4}\right)^2
\bigg].
\end{split}
\label{eq:h-bessel}
\end{equation}
Felder and Hoppe proved the required Gram nondegeneracy for
\(d=1,2,3\) and formulated it as a conjecture for general \(d\).
The quartic case is the first one not covered by their argument.

Our result is the following.

\begin{theorem}\label{thm:main}
For every \(a>0\), the sequence determined by
\eqref{eq:quartic-recurrence}--\eqref{eq:h-bessel} satisfies
\[
v_n(a)>0\qquad\text{for every }n\geq0.
\]
\end{theorem}

The main point is an integral sign mechanism that applies uniformly to
all indices.  The Bessel expression \eqref{eq:h-bessel} has a positive
Laplace representation.  Adding independent source variables for
\(z^4\) and \(\bar z^4\) produces four positive integral kernels, one
for each residue class modulo \(4\).  Their mixed derivative matrices
are reduced by triangular row operations to
\[
\left[
\int_0^\infty s^p
\left(\frac{s^2}{1+s^2}\right)^q\dd\mu(s)
\right]_{p,q=0}^{k-1}.
\]
Both \(s\) and \(s^2/(1+s^2)\) are strictly increasing.  Andr\'eief's
identity therefore expresses the determinant as an integral of a
product of two Vandermonde determinants with a fixed positive sign.
This proves the required sign pattern for every finite block at once.

The paper is organized as follows.  Section~\ref{sec:gram} recalls the
quartic Hermitian form and its relation to the recurrence.
Section~\ref{sec:sign-lemma} proves the abstract sign lemma.
Section~\ref{sec:blocks} derives the four positive kernels and applies
the lemma to the four symmetry blocks.  Section~\ref{sec:completion}
then identifies the positive norm ratios with the solution of
\eqref{eq:quartic-recurrence}.

\section{The quartic Gram form}\label{sec:gram}

Fix \(a>0\) and write
\begin{equation}
\rho_a(z)=
\exp\left(-a|z|^2+\frac{\ii}{4}(z^4+\bar z^4)\right).
\label{eq:density}
\end{equation}
The Gaussian factor makes all polynomial moments absolutely
convergent.  Let \(\eta=\e^{\pi\ii/4}\).  Decompose the polynomial ring
into \(C_4\)-weight spaces
\[
\C[z]=\bigoplus_{r=0}^3\mathcal P_r,
\qquad
\mathcal P_r=\{f\in\C[z]:f(\ii z)=\ii^r f(z)\}.
\]
For \(f\in\mathcal P_r\), set
\[
(\sigma f)(z)=(-\eta)^{-r}f(\eta z),
\]
and define the Hermitian form
\begin{equation}
(f,g)_a
=\int_{\C}f(z)\overline{\sigma g(z)}\rho_a(z)\dd A(z).
\label{eq:pairing}
\end{equation}
The density is invariant under \(z\mapsto\ii z\), so distinct weight
spaces are orthogonal.  If \(m=r+4q\), then
\begin{equation}
\sigma z^m=(-\eta)^{-r}\eta^m z^m=(-1)^{r+q}z^m.
\label{eq:twist-sign}
\end{equation}
Thus the Gram matrix on \(\mathcal P_r\), in the ordered monomial basis
\(z^r,z^{r+4},\ldots\), has entries
\begin{equation}
G_{pq}^{(r)}
=(-1)^{r+q}\int_{\C}
z^{r+4p}\bar z^{r+4q}\rho_a(z)\dd A(z),
\qquad p,q\geq0.
\label{eq:gram-block}
\end{equation}

Assume temporarily that the leading principal minors in all four
blocks are nonzero.  Gram--Schmidt then gives monic polynomials
\(P_n(z)=z^n+\cdots\), and symmetry implies
\(P_n\in\mathcal P_{n\bmod4}\).  Define
\begin{equation}
\tau(\ell)=
\begin{cases}
-1,&\ell\equiv3\pmod4,\\
1,&\ell\not\equiv3\pmod4,
\end{cases}
\qquad
\varrho(n)=(-1)^n\prod_{\ell=0}^{n-1}\tau(\ell)
=(-1)^{n+\lfloor n/4\rfloor}.
\label{eq:rho-sign}
\end{equation}
The signed norms \(h_n\) are fixed by
\begin{equation}
(P_n,P_n)_a=\varrho(n)h_n.
\label{eq:signed-norm}
\end{equation}
Section~\ref{sec:blocks} will prove simultaneously that all the
required minors are nonzero and that \(h_n>0\).

\section{A sign-regularity lemma}\label{sec:sign-lemma}

We isolate the integral argument that drives the proof.  Let
\(y>0\), \(\alpha\geq0\), and \(\beta\geq0\), and define
\begin{equation}
\Phi_{\alpha,\beta}(w;y)
=\int_0^\infty
\e^{-ys}s^{\alpha-1/2}(1+ws^2)^{-\beta-1/2}\dd s,
\qquad w>0.
\label{eq:Phi}
\end{equation}
For \(p,q\geq0\), put
\begin{equation}
D_{pq}^{\alpha,\beta}
=\left.
\partial_u^p\partial_v^q
\Phi_{\alpha,\beta}(uv;y)
\right|_{u=v=1}.
\label{eq:D-def}
\end{equation}

\begin{lemma}[strict sign regularity]\label{lem:sign-regularity}
For every \(k\geq1\),
\begin{equation}
\sgn\det
\left[D_{pq}^{\alpha,\beta}\right]_{p,q=0}^{k-1}
=(-1)^{k(k-1)/2}.
\label{eq:D-sign}
\end{equation}
\end{lemma}

\begin{proof}
Set
\[
b=\beta+\frac12,
\qquad
\gamma=\frac{\alpha}{2}+\frac14,
\qquad
L(y)=\Phi_{\alpha,\beta}(1;y).
\]
The substitution \(t=\sqrt w\,s\) in \eqref{eq:Phi} gives
\begin{equation}
\Phi_{\alpha,\beta}(w;y)
=w^{-\gamma}L\left(\frac{y}{\sqrt w}\right).
\label{eq:Phi-scaling}
\end{equation}
Let \(\vartheta=w\partial_w\), and use the falling-factorial notation
\[
(\vartheta)\fall{p}
=\vartheta(\vartheta-1)\cdots(\vartheta-p+1).
\]
Under \eqref{eq:Phi-scaling}, \(\vartheta\) acts on \(L\) as
\[
\mathcal A=-\gamma-\frac y2\partial_y.
\]
Since \(w=uv\),
\[
\left.
\partial_u^p\partial_v^q F(uv)
\right|_{u=v=1}
=\left.
(\vartheta)\fall{p}(\vartheta)\fall{q}F(w)
\right|_{w=1}.
\]
It follows that
\begin{equation}
D_{pq}^{\alpha,\beta}
=(\mathcal A)\fall{p}(\mathcal A)\fall{q}L(y).
\label{eq:D-operator}
\end{equation}

Define
\[
f(s)=\frac{s^2}{1+s^2}
\]
and the bi-moments
\begin{equation}
B_{pq}
=\int_0^\infty
s^p f(s)^q
\e^{-ys}s^{\alpha-1/2}(1+s^2)^{-\beta-1/2}\dd s.
\label{eq:B-moments}
\end{equation}
Direct differentiation of \eqref{eq:Phi} with respect to \(w\) gives
\begin{equation}
(\mathcal A)\fall{q}L(y)
=(-1)^q(b)_qB_{0q},
\label{eq:A-on-L}
\end{equation}
where \((b)_q\) is the rising Pochhammer symbol.  On the other hand,
\begin{equation}
(\mathcal A)\fall{p}
=\sum_{i=0}^p c_{pi}(y)\partial_y^i,
\qquad
c_{pp}(y)=\left(-\frac y2\right)^p.
\label{eq:A-expansion}
\end{equation}
Since \(\partial_y^iB_{0q}=(-1)^iB_{iq}\),
\eqref{eq:D-operator}--\eqref{eq:A-expansion} give a factorization
\begin{equation}
D=RBC,
\label{eq:triangular-factorization}
\end{equation}
where \(R\) is lower triangular with diagonal
\[
R_{pp}=\left(\frac y2\right)^p>0,
\]
and \(C\) is diagonal with
\[
C_{qq}=(-1)^q(b)_q.
\]

It remains to determine the sign of \(\det B\).  Andr\'eief's
integration identity \cite{Andreief1886,Forrester2018} gives
\begin{equation}
\begin{split}
\det[B_{pq}]_{p,q=0}^{k-1}
=\frac1{k!}\int_{(0,\infty)^k}
&\det[s_j^p]_{p=0,\ldots,k-1}^{j=1,\ldots,k}\\
{}\times&
\det[f(s_j)^q]_{q=0,\ldots,k-1}^{j=1,\ldots,k}
\prod_{j=1}^k\dd\mu(s_j),
\end{split}
\label{eq:andreief}
\end{equation}
where
\[
\dd\mu(s)
=\e^{-ys}s^{\alpha-1/2}
(1+s^2)^{-\beta-1/2}\dd s.
\]
The two determinants in \eqref{eq:andreief} are Vandermonde
determinants.  Because both \(s\) and
\(f(s)=s^2/(1+s^2)\) are strictly increasing on \((0,\infty)\),
their product is nonnegative and is strictly positive whenever the
\(s_j\)'s are distinct.  Hence \(\det B>0\).

The determinant of \(R\) is positive, whereas
\[
\sgn\det C=(-1)^{0+1+\cdots+(k-1)}
=(-1)^{k(k-1)/2}.
\]
The factorization \eqref{eq:triangular-factorization} proves
\eqref{eq:D-sign}.
\end{proof}

\section{The four positive kernels}\label{sec:blocks}

We first turn the Bessel initial function into a positive integral.
For \(x>0\), the connection formula
\[
J_{-1/4}(x)
=\frac{J_{1/4}(x)-Y_{1/4}(x)}{\sqrt2}
\]
reduces \eqref{eq:h-bessel} to
\begin{equation}
h(a)=\frac{\pi^2a}{8}
\left[
J_{1/4}\left(\frac{a^2}{4}\right)^2
+Y_{1/4}\left(\frac{a^2}{4}\right)^2
\right].
\label{eq:h-amplitude}
\end{equation}
The Dixon--Ferrar formula \cite{DixonFerrar1933,GradshteynRyzhik2015}
states that, for \(|\operatorname{Re}\nu|<\frac12\),
\begin{equation}
J_\nu(x)^2+Y_\nu(x)^2
=\frac{8\cos(\pi\nu)}{\pi^2}
\int_0^\infty K_{2\nu}(2x\sinh t)\dd t.
\label{eq:dixon-ferrar}
\end{equation}
Using
\[
K_{1/2}(u)=\sqrt{\frac{\pi}{2u}}\e^{-u}
\]
in \eqref{eq:h-amplitude}, and then setting \(s=\sinh t\), gives
\begin{equation}
h(a)=\sqrt{\frac{\pi}{2}}\int_0^\infty
\frac{\e^{-a^2s/2}}{\sqrt{s(1+s^2)}}\dd s.
\label{eq:h-positive}
\end{equation}

Set
\[
y=\frac{a^2}{2},
\qquad
C=\sqrt{\frac{\pi}{2}}.
\]
We now introduce an angular source.  For \(u,v>0\), define
\begin{equation}
\mathcal Z(a;u,v)
=\pi\int_0^\infty
\e^{-at}J_0\left(\frac{t^2\sqrt{uv}}2\right)\dd t.
\label{eq:source}
\end{equation}
This is the angularly integrated two-source expression.  Indeed,
termwise angular integration gives
\begin{equation}
\int_0^{2\pi}
\exp\left(
\frac{\ii R^4}{4}
(u\e^{4\ii\theta}+v\e^{-4\ii\theta})
\right)\dd\theta
=2\pi J_0\left(\frac{R^4\sqrt{uv}}2\right),
\label{eq:angular-source}
\end{equation}
and \(t=R^2\) gives \eqref{eq:source}.  Only derivatives at
\(u=v=1\) will be used, where all resulting planar moment integrals are
absolutely convergent.

At \(u=v=1\), \(\mathcal Z(a;1,1)=h(a)\).  Scaling \(t\) in
\eqref{eq:source} and using \eqref{eq:h-positive} yield
\begin{equation}
\mathcal Z(a;u,v)
=C\int_0^\infty
\frac{\e^{-ys}}{\sqrt{s(1+uvs^2)}}\dd s
=C\Phi_{0,0}(uv;y).
\label{eq:source-positive}
\end{equation}
For \(r=0,1,2,3\), put
\[
\mathcal Z_r=(-\partial_a)^r\mathcal Z.
\]

\begin{lemma}\label{lem:four-kernels}
The four radial source functions are
\begin{equation}
\begin{array}{c|c}
r&\mathcal Z_r(a;u,v)\\ \hline
0&C\Phi_{0,0}(uv;y),\\[1mm]
1&aC\Phi_{1,0}(uv;y),\\[1mm]
2&2C\Phi_{1,1}(uv;y),\\[1mm]
3&2aC\Phi_{2,1}(uv;y).
\end{array}
\label{eq:four-kernels}
\end{equation}
\end{lemma}

\begin{proof}
The first two rows follow by differentiating
\eqref{eq:source-positive}.  For the third row, integration by parts
gives
\begin{equation}
\begin{split}
y\int_0^\infty
\frac{\e^{-ys}s^{3/2}}{\sqrt{1+uvs^2}}\dd s
={}&\frac12\int_0^\infty
\frac{\e^{-ys}s^{1/2}}{\sqrt{1+uvs^2}}\dd s\\
&+\int_0^\infty
\frac{\e^{-ys}s^{1/2}}{(1+uvs^2)^{3/2}}\dd s.
\end{split}
\label{eq:ibp-kernel}
\end{equation}
The boundary terms vanish at \(0\) and \(\infty\).  Since
\[
(-\partial_a)^2
=\partial_y+2y\partial_y^2,
\]
equation \eqref{eq:ibp-kernel} gives
\(\mathcal Z_2=2C\Phi_{1,1}\).  One more differentiation gives
\(\mathcal Z_3=2aC\Phi_{2,1}\).
\end{proof}

For \(k\geq1\), set
\begin{equation}
\Delta_k^{(r)}
=\det[G_{pq}^{(r)}]_{p,q=0}^{k-1},
\qquad
\Delta_0^{(r)}=1.
\label{eq:Delta}
\end{equation}

\begin{proposition}\label{prop:block-signs}
For every \(r\in\{0,1,2,3\}\) and \(k\geq1\),
\begin{equation}
\sgn\Delta_k^{(r)}
=(-1)^{rk+k(k-1)/2}.
\label{eq:block-sign}
\end{equation}
\end{proposition}

\begin{proof}
Let
\[
D_{pq}^{(r)}
=\left.
\partial_u^p\partial_v^q\mathcal Z_r(a;u,v)
\right|_{u=v=1}.
\]
Every scalar multiplier in \eqref{eq:four-kernels} is positive.
Lemma~\ref{lem:sign-regularity}, applied to the four parameter pairs
\[
(\alpha,\beta)=(0,0),(1,0),(1,1),(2,1),
\]
therefore gives
\begin{equation}
\sgn\det[D_{pq}^{(r)}]_{p,q=0}^{k-1}
=(-1)^{k(k-1)/2}.
\label{eq:Dr-sign}
\end{equation}

Differentiating \eqref{eq:angular-source} at \(u=v=1\) shows that
\[
D_{pq}^{(r)}
=\left(\frac{\ii}{4}\right)^{p+q}
\int_{\C}
z^{r+4p}\bar z^{r+4q}\rho_a(z)\dd A(z).
\]
Combining this with \eqref{eq:gram-block} gives
\begin{equation}
G_{pq}^{(r)}
=(-1)^{r+q}\left(\frac4\ii\right)^{p+q}
D_{pq}^{(r)}.
\label{eq:G-D}
\end{equation}
Let \(S_k=k(k-1)/2\).  Factoring the row and column phases in
\eqref{eq:G-D} yields
\begin{equation}
\Delta_k^{(r)}
=(-1)^{rk}16^{S_k}
\det[D_{pq}^{(r)}]_{p,q=0}^{k-1}.
\label{eq:Delta-D}
\end{equation}
Now \eqref{eq:block-sign} follows from \eqref{eq:Dr-sign}.
\end{proof}

\begin{corollary}\label{cor:positive-norms}
All quartic Gram matrices are nondegenerate, and the signed norms in
\eqref{eq:signed-norm} satisfy
\[
h_n>0\qquad(n\geq0).
\]
\end{corollary}

\begin{proof}
The block sign formula shows first that every
\(\Delta_k^{(r)}\) is nonzero.  Thus all the monic orthogonal
polynomials exist.  If \(n=r+4j\), Gram--Schmidt within
\(\mathcal P_r\) gives
\begin{equation}
(P_n,P_n)_a
=\frac{\Delta_{j+1}^{(r)}}{\Delta_j^{(r)}}.
\label{eq:norm-delta-ratio}
\end{equation}
Taking the quotient of the signs in \eqref{eq:block-sign} gives
\[
\sgn(P_n,P_n)_a=(-1)^{r+j}.
\]
But \eqref{eq:rho-sign} gives
\(\varrho(r+4j)=(-1)^{r+j}\).  Hence
\(h_n=\varrho(n)(P_n,P_n)_a>0\).
\end{proof}

\section{Identification with the quartic recurrence}
\label{sec:completion}

It remains to verify that the positive norm ratios give precisely the
sequence in Theorem~\ref{thm:main}.  Define
\begin{equation}
V_n=\frac{h_{n+1}}{h_n},
\qquad
v_n=\frac{V_n}{a}.
\label{eq:ratios}
\end{equation}
Corollary~\ref{cor:positive-norms} already gives
\begin{equation}
v_n>0\qquad(n\geq0).
\label{eq:ratio-positive}
\end{equation}

For clarity, we include the short recurrence calculation from the
quartic Gram form.  Differentiating \eqref{eq:density} and integrating
by parts in \eqref{eq:pairing} give
\begin{equation}
\big((\partial_z+\ii z^3)f,g\big)_a
=-\tau(\ell)a(f,zg)_a
\label{eq:ibp-pairing}
\end{equation}
when \(g\) has \(C_4\)-weight \(\ell\).  The Gaussian decay removes
the boundary term.  Weight, degree, and orthogonality imply
\begin{equation}
\begin{split}
zP_n
&=P_{n+1}+\frac{\ii}{a}\frac{h_n}{h_{n-3}}P_{n-3},\\
(\partial_z+\ii z^3)P_n
&=\ii P_{n+3}+a\frac{h_n}{h_{n-1}}P_{n-1},
\end{split}
\label{eq:polynomial-recurrences}
\end{equation}
where terms with negative indices are omitted.  To see the coefficients,
pair the first expansion with \(P_{n-3}\), the second with
\(P_{n-1}\), and use
\[
\varrho(n)=\tau(n)\varrho(n-3),
\qquad
\varrho(n)=-\tau(n-1)\varrho(n-1).
\]

The coefficient of \(P_{n-1}\) in \(z^3P_n\) receives three
contributions: the downward step in the first formula of
\eqref{eq:polynomial-recurrences} may occur after zero, one, or two
upward steps.  Hence that coefficient is
\[
\frac{\ii}{a}V_{n-1}
\left(
V_{n-2}V_{n-3}
+V_nV_{n-2}
+V_{n+1}V_n
\right).
\]
Subtracting \(\ii z^3P_n\) from the second formula in
\eqref{eq:polynomial-recurrences}, and using that the coefficient of
\(P_{n-1}\) in \(\partial_zP_n\) is \(n\), gives
\begin{equation}
V_m\left[
1+\frac1{a^2}
\left(
V_{m-1}V_{m-2}
+V_{m+1}V_{m-1}
+V_{m+2}V_{m+1}
\right)
\right]
=\frac{m+1}{a}.
\label{eq:V-recurrence}
\end{equation}
With \(V_m=av_m\), this is
\eqref{eq:quartic-recurrence}.

Finally, \(P_0=1\), \(P_1=z\), and \(P_2=z^2\).  Differentiating
\(h(a)=(1,1)_a\) gives
\[
h_0=h,\qquad h_1=-h',\qquad h_2=h''.
\]
Therefore \eqref{eq:ratios} gives exactly
\eqref{eq:initial-values}.  Conversely, the forward form of
\eqref{eq:quartic-recurrence} determines \(v_{m+2}\) from the preceding
terms whenever \(v_mv_{m+1}\neq0\).  The positive sequence
\eqref{eq:ratios} has this property, so induction identifies it with
the sequence determined by \eqref{eq:initial-values}.  Together with
\eqref{eq:ratio-positive}, this proves Theorem~\ref{thm:main}.

\begin{remark}
The argument proves only the positivity assertion required in
Theorem~\ref{thm:main}.  It does not use or claim a classification of
all positive solutions, a closed evaluation of the Gram determinants,
or asymptotic formulas for \(v_n\).
\end{remark}

\bibliographystyle{amsplain}
\bibliography{references}

\end{document}
